% !TEX program = pdflatex
\documentclass[11pt,a4paper]{article}

\usepackage[margin=1in]{geometry}
\usepackage{amsmath,amssymb,amsthm,mathtools}
\usepackage[T1]{fontenc}
\usepackage{lmodern}
\usepackage{microtype}
\usepackage{booktabs}
\usepackage{array}
\usepackage{enumitem}
\usepackage{hyperref}
\usepackage{xcolor}
\usepackage{longtable}
\usepackage{siunitx}
\usepackage{listings}
\usepackage{fancyhdr}
\usepackage{lastpage}

\hypersetup{
  colorlinks=true,
  linkcolor=blue!50!black,
  citecolor=blue!50!black,
  urlcolor=blue!50!black,
  pdftitle={Arithmetic-Sequenced Polynomial Coordinates for Semiprime Factorization},
  pdfauthor={IamLupo}
}

\newtheorem{theorem}{Theorem}[section]
\newtheorem{proposition}[theorem]{Proposition}
\newtheorem{lemma}[theorem]{Lemma}
\newtheorem{corollary}[theorem]{Corollary}
\newtheorem{definition}[theorem]{Definition}
\newtheorem{remark}[theorem]{Remark}
\newtheorem{conjecture}[theorem]{Conjecture}

\newcommand{\N}{N}
\newcommand{\R}{\mathbb{R}}
\newcommand{\Z}{\mathbb{Z}}
\newcommand{\Fp}{\mathbb{F}}
\newcommand{\vp}{v_2}
\newcommand{\floor}[1]{\left\lfloor #1 \right\rfloor}
\newcommand{\ceil}[1]{\left\lceil #1 \right\rceil}
\newcommand{\dd}{\,\mathrm d}

\setlength{\headheight}{14pt}
\pagestyle{fancy}
\fancyhf{}
\lhead{Arithmetic-Sequenced Polynomial Coordinates}
\rhead{\thepage/\pageref{LastPage}}
\renewcommand{\headrulewidth}{0.4pt}

\lstset{
  basicstyle=\ttfamily\small,
  columns=fullflexible,
  frame=single,
  breaklines=true,
  showstringspaces=false
}

\title{\textbf{Arithmetic-Sequence Coordinates for Semiprime Factorization}\\
\large From Binomial Identities to an Exact Divisor Parameterization, 2-Adic Structure, and Quadratic-Residue Sieving}
\author{Experimental Mathematics Notes}
\date{September 2026}

\begin{document}
\maketitle
\begin{abstract}
We investigate a coordinate system for factoring an odd semiprime \(N=pq\), obtained by combining an arithmetic-sequence/binomial transform with a polynomial relation between a near-square coordinate and a sum-gap coordinate. The construction begins with the binomial identities
\[
\binom{x}{y}2^{x-y}=[t^y](2+t)^x,\qquad
\sum_y\binom{x}{y}2^{x-y}=3^x,
\]
and their first-moment analogue. These identities convert a quadratic polynomial \(P(x,y)\) into a one-variable quadratic \(H(x)\), while a second use of the same transform on the diagonal \(x+y=z\) linearizes the polynomial exactly. The resulting diagonal parameter \(r=s+z\) is not merely correlated with a factor: \(r\mid N\) is equivalent to the existence of an integer zero on the corresponding diagonal, and the finite difference of the transformed sequence is exactly \(r\).

The same construction also admits several equivalent descriptions. The \(y\)-coordinate reproduces Fermat factorization exactly; the discriminant of \(P\) is the squared factor gap \((q-p)^2\). An alternative parametrization by \(a=s-p\) and \(b=q-s\) converts the problem into the classical relation \(ab=s(y+1)-D\). Reduction modulo powers of two yields an exact parity condition, an involution on the modular roots, and a connection to square roots of \(N\) modulo \(2^k\). These 2-adic structures are exact but, in the experiments considered here, do not by themselves expose the factors. Finally, a quadratic-residue sieve on the Fermat coordinate \(y\) can reduce the candidate set by roughly three orders of magnitude, although the observed rank of the true \(y\) is quantitatively consistent with ordinary sieve density rather than an additional factor-localizing principle.

The central outcome is therefore a clean algebraic framework with several exact identities, a useful high-strength modular sieve, and a clear separation between identities that merely reparameterize Fermat factorization and mechanisms that would constitute a genuinely new factorization method.
\end{abstract}

\tableofcontents
\newpage

\section{Arithmetic sequences and the binomial identities}
\label{sec:arithmetic}

The starting point is the weighted binomial sequence
\begin{equation}
W(x,y)=\binom{x}{y}2^{x-y},
\qquad 0\le y\le x,
\label{eq:W}
\end{equation}
which is naturally generated by the arithmetic shift of the binomial theorem.

\subsection{Generating-function interpretation}

Since
\begin{equation}
(2+t)^x=\sum_{y=0}^{x}\binom{x}{y}2^{x-y}t^y,
\end{equation}
we obtain immediately
\begin{equation}
\boxed{\binom{x}{y}2^{x-y}=[t^y](2+t)^x.}
\label{eq:coefficient}
\end{equation}
Consequently,
\begin{equation}
\sum_{y=0}^{x}\binom{x}{y}2^{x-y}=3^x.
\label{eq:sum0}
\end{equation}
Differentiating with respect to the auxiliary variable gives the first moment
\begin{equation}
\sum_{y=0}^{x}y\binom{x}{y}2^{x-y}=x3^{x-1}.
\label{eq:sum1}
\end{equation}
The second moment, used later as a consistency check, is
\begin{equation}
\sum_{y=0}^{x}y(y-1)\binom{x}{y}2^{x-y}=x(x-1)3^{x-2},
\end{equation}
and hence
\begin{equation}
\sum_{y=0}^{x}y^2\binom{x}{y}2^{x-y}
=x(x-1)3^{x-2}+x3^{x-1}.
\label{eq:sum2}
\end{equation}

These formulas imply that any polynomial of degree at most two in the summation variable can be evaluated exactly in closed form under the weight \(W(x,y)\). This is the arithmetic-sequence entry point of the construction.

\subsection{A general transform identity}

For a polynomial
\[
Q(x,y)=A(x)+B(x)y+C(x)y^2,
\]
its weighted transform
\[
\mathcal T[Q](x)=\sum_{y=0}^{x}Q(x,y)\binom{x}{y}2^{x-y}
\]
is
\begin{equation}
\mathcal T[Q](x)
=3^{x-2}\left(9A(x)+3xB(x)+x(x+2)C(x)\right).
\label{eq:generaltransform}
\end{equation}
This is simply the substitution of \eqref{eq:sum0}--\eqref{eq:sum2}. In particular, no approximation is involved.

\section{The semiprime coordinate system}
\label{sec:coordinates}

Let
\[
N=pq
\]
be an odd semiprime with distinct odd primes \(p<q\). Define
\begin{equation}
s=\floor{\sqrt N},
\qquad
D=N-s^2,
\label{eq:sD}
\end{equation}
and introduce the coordinates
\begin{equation}
\boxed{x=s+1-p},
\qquad
\boxed{y=p+q-2s-1}.
\label{eq:xytrue}
\end{equation}
Equivalently,
\begin{equation}
\boxed{p=s+1-x},
\qquad
\boxed{q=s+x+y}.
\label{eq:pqxy}
\end{equation}

The quantity \(D=N-s^2\) measures the displacement of \(N\) above the nearest lower square. The coordinate \(x\) measures how far the smaller factor lies below \(s+1\), while \(y\) measures the offset of the factor sum from \(2s+1\).

\subsection{The polynomial}

Substitution of \eqref{eq:pqxy} into \(N-pq\) gives
\begin{equation}
\boxed{
P(x,y)=x^2+(y-1)x+(D-s)-(s+1)y.
}
\label{eq:P}
\end{equation}
The crucial identity is
\begin{equation}
\boxed{
P(x,y)=N-(s+1-x)(s+x+y).
}
\label{eq:Pfactor}
\end{equation}
At the true coordinates \((x_p,y)\), we therefore have
\begin{equation}
P(x_p,y)=0.
\end{equation}
The same value of \(y\) has a second root in \(x\), corresponding to exchanging the factors. Since the coefficient of \(x\) is \(y-1\), Vieta gives
\begin{equation}
\boxed{x_p+x_q=1-y.}
\label{eq:vieta}\end{equation}
Here
\[
x_q=s+1-q<0.
\]
Moreover,
\begin{equation}
\boxed{x_p-x_q=q-p.}
\label{eq:gapx}
\end{equation}

\section{The diagonal: an exact divisor parameterization}
\label{sec:diagonal}

The key transformation is to fix the diagonal variable
\begin{equation}
z=x+y.
\end{equation}
Then \(y=z-x\), and substituting into \eqref{eq:P} yields
\begin{align}
P(x,z-x)
&=x^2+(z-x-1)x+(D-s)-(s+1)(z-x)\\
&=(z+s)x+D-s-(s+1)z.
\end{align}
Thus
\begin{equation}
\boxed{P(x,z-x)=(z+s)x+D-s-(s+1)z.}
\label{eq:diaglinear}
\end{equation}
The quadratic term has disappeared completely.

Define
\begin{equation}
r=s+z.
\end{equation}
Since \(D=N-s^2\), the diagonal expression can be rewritten as
\begin{equation}
\boxed{
P(x,z-x)=rx+N-r(s+1).
}
\label{eq:diagN}
\end{equation}
Equivalently,
\begin{equation}
\boxed{
P(x,z-x)=N-r(s+1-x).
}
\label{eq:diagfactor}
\end{equation}

\begin{theorem}[Diagonal divisibility equivalence]
For integers \(x,z\), let \(r=s+z\). Then
\begin{equation}
P(x,z-x)=0
\end{equation}
if and only if
\begin{equation}
r\mid N.
\end{equation}
More precisely, when \(r\mid N\), the unique corresponding integer \(x\) is
\begin{equation}
\boxed{x=s+1-\frac Nr.}
\label{eq:xfromr}
\end{equation}
\end{theorem}

\begin{proof}
From \eqref{eq:diagfactor},
\[
P(x,z-x)=0
\iff
N=r(s+1-x).
\]
Thus \(r\mid N\), with quotient \(N/r=s+1-x\). Conversely, any divisor \(r\mid N\) inserted into \eqref{eq:xfromr} gives an integer zero. \qed
\end{proof}

The diagonal coordinate therefore converts the original two-variable zero problem into a divisor parameterization. At \(r=p\) or \(r=q\), the corresponding zero is one of the two true roots.

\subsection{Exact finite difference}

Along a fixed diagonal \(z\), incrementing \(x\) by one necessarily decrements \(y\) by one. From \eqref{eq:diaglinear},
\begin{equation}
\boxed{
P(x+1,z-x-1)-P(x,z-x)=z+s=r.
}
\label{eq:slope}
\end{equation}
Hence the diagonal slope is literally the divisor parameter. This fact will reappear in the binomial transform.

\section{The binomial transform of the original polynomial}
\label{sec:H}

Apply the arithmetic-sequence transform to \(P(x,y)\):
\begin{equation}
T(x)=\sum_{y=0}^{x}P(x,y)\binom{x}{y}2^{x-y}.
\label{eq:T}
\end{equation}
Using \eqref{eq:sum0} and \eqref{eq:sum1}, one obtains
\begin{equation}
\boxed{
T(x)=3^{x-1}H(x),
}
\label{eq:T-H}
\end{equation}
where
\begin{equation}
\boxed{
H(x)=4x^2-(s+4)x+3D-3s.
}
\label{eq:H}
\end{equation}
An equivalent and useful shifted form is
\begin{equation}
\boxed{
H(x)=3N+(x-s-1)(4x+3s).
}
\label{eq:Hshift}
\end{equation}
Therefore, modulo any divisor \(r\mid N\),
\begin{equation}
\boxed{
H(x)\equiv(x-s-1)(4x+3s)\pmod r.
}
\label{eq:Hmod}
\end{equation}
At the true factor coordinate \(x_p=s+1-p\), one obtains
\begin{equation}
\boxed{
H(x_p)=p(4p+3q-7s-4).
}
\label{eq:Hxtrue}
\end{equation}

\subsection{Why ordinary features of \(H\) do not locate the factor}

The polynomial \(H\) is exact, but several natural attempts to use its real geometry were tested experimentally. In a 5000-case study, all cases had real roots. Nevertheless, the true coordinate \(x_p\) was far from the roots and the vertex on average, local minima never identified \(x_p\), and simple root/vertex/minimum predictors produced zero correct predictions out of 5000 cases. The exact identity \eqref{eq:Hxtrue} was verified in every test.

The conclusion is important: the transform compresses the two-variable arithmetic into a clean quadratic, but this quadratic's ordinary real geometry does not encode the smaller factor location in a useful way.

\section{The binomial transform on a diagonal}
\label{sec:diagonaltransform}

The same transform behaves much more rigidly after diagonalization. Fix \(r=s+z\), so that
\[
P(x,z-x)=rx+N-r(s+1).
\]
Define
\begin{equation}
B_n(r)=\sum_{x=0}^{n}P(x,z-x)\binom{n}{x}2^{n-x}.
\label{eq:Bn}
\end{equation}
Since the diagonal expression is linear in \(x\), the identities \eqref{eq:sum0} and \eqref{eq:sum1} give
\begin{align}
B_n(r)
&=r\,n3^{n-1}+\bigl(N-r(s+1)\bigr)3^n\\
&=3^{n-1}\left[3N+r\bigl(n-3(s+1)\bigr)\right].
\end{align}
Hence
\begin{equation}
\boxed{
B_n(r)=3^{n-1}\left[3N+r(n-3(s+1))\right].
}
\label{eq:Bnclosed}
\end{equation}
Define the normalized quantity
\begin{equation}
C_n(r)=\frac{B_n(r)-3^nN}{3^{n-1}}.
\end{equation}
Then
\begin{equation}
\boxed{C_n(r)=r\,[n-3(s+1)]}
\label{eq:Cn}
\end{equation}
and therefore
\begin{equation}
\boxed{C_{n+1}(r)-C_n(r)=r.}
\label{eq:Cdifference}
\end{equation}

Thus the arithmetic-sequence transform does not merely preserve the diagonal slope: after normalization, its first difference is exactly the divisor parameter. Experiments recovered known factors from this transformed difference with 100\% success whenever the diagonal radius \(r\) was already supplied.

However, scanning candidate radii produced the same search problem as directly testing
\[
\gcd(r,N).
\]
The transform adds no observed directional information for selecting the correct radius. This distinction is central: \emph{encoding a known divisor more elegantly is not equivalent to discovering the divisor}.

\section{The \(y\)-coordinate and Fermat factorization}
\label{sec:fermat}

The second coordinate has an equally exact classical interpretation. Treat \(P(x,y)\) as a quadratic in \(x\). Its discriminant is
\begin{align}
\Delta(y)
&=(y-1)^2-4\bigl[(D-s)-(s+1)y\bigr]\\
&=(2s+1+y)^2-4N.
\end{align}
Therefore
\begin{equation}
\boxed{
\Delta(y)=(2s+1+y)^2-4N.
}
\label{eq:disc}
\end{equation}
At the true coordinate \(y=p+q-2s-1\),
\begin{equation}
2s+1+y=p+q,
\end{equation}
so
\begin{equation}
\boxed{
\Delta(y)=(p+q)^2-4pq=(q-p)^2.
}
\label{eq:gap}
\end{equation}

Thus the exact factor test is simply:
\begin{equation}
(2s+1+y)^2-4N
\end{equation}
is a perfect square. This is precisely Fermat factorization in shifted coordinates.

If
\[
A=s+\ceil{y/2},
\]
then the classical Fermat quantity is
\begin{equation}
A^2-N=\frac{(q-p)^2}{4}
\end{equation}
at the correct \(y\). Experiments showed zero mismatches between the \(y\)-search and ordinary Fermat iteration, apart from the expected change of coordinate origin.

The consequence is both positive and negative. Positive: the polynomial framework recovers Fermat factorization exactly and exposes its algebraic structure. Negative: searching \(y\) by itself is not a new factorization algorithm.

\section{The alternative \(a,b\) coordinates}
\label{sec:ab}

Define
\begin{equation}
a=s-p\ge0,
\qquad
b=q-s>0.
\label{eq:ab}
\end{equation}
Then
\begin{equation}
p=s-a,
\qquad q=s+b.
\end{equation}
The gap relation becomes
\begin{equation}
q-p=a+b.
\end{equation}
For the \(y\)-coordinate,
\begin{equation}
\boxed{y=b-a-1,}
\label{eq:yab}
\end{equation}
so
\begin{equation}
\boxed{b-a=y+1.}
\label{eq:diffab}
\end{equation}
Expanding \(N=(s-a)(s+b)\) gives
\begin{align}
D
&=N-s^2\\
&=s(b-a)-ab\\
&=s(y+1)-ab.
\end{align}
Hence
\begin{equation}
\boxed{ab=s(y+1)-D.}
\label{eq:abproduct}
\end{equation}

Equations \eqref{eq:diffab} and \eqref{eq:abproduct} determine \(a\) and \(b\) once \(y\) is known. Indeed, letting \(d=y+1\), one has
\[
b-a=d,
\qquad ab=s d-D,
\]
and therefore
\begin{equation}
(a+b)^2=d^2+4(sd-D),
\end{equation}
which is again the discriminant/gap relation in another notation.

Experiments verified all these identities over thousands of random semiprimes. Scanning \(y\) and requiring the resulting quadratic to have integral roots yields only the classical square test. The \(ab\) reformulation is therefore an exact restatement rather than an independent shortcut.

\section{Modular analysis modulo powers of two}
\label{sec:mod2}

The modular structure is considerably richer. Let \(N\) be odd and work modulo \(2^k\). Write
\begin{equation}
P(x,y)=A(x)+B(x)y,
\end{equation}
where
\begin{align}
A(x)&=x^2-x+D-s,\\
B(x)&=x-s-1.
\end{align}
A congruence
\[
P(x,y)\equiv0\pmod{2^k}
\]
has a solution in \(y\) exactly when
\[
\gcd(B(x),2^k)\mid A(x).
\]
But
\begin{equation}
A(x)=(x-s)(x-s-1)+D,
\end{equation}
and modulo any common divisor of \(B(x)=x-s-1\) and \(2^k\), the term involving \(x-s-1\) vanishes while \(A(x)\equiv N\). Since \(N\) is odd, the divisibility condition forces \(B(x)\) to be odd.

Thus
\begin{equation}
\boxed{
P(x,y)\equiv0\pmod{2^k}
\quad\text{is solvable in }y
\iff x\equiv s\pmod2.
}
\label{eq:parity}
\end{equation}
When this parity condition holds, \(B(x)\) is invertible modulo \(2^k\), and the corresponding solution is unique:
\begin{equation}
\boxed{
 y\equiv-\bigl(x^2-x+D-s\bigr)(x-s-1)^{-1}\pmod{2^k}.
}
\label{eq:y2adic}
\end{equation}
Therefore exactly \(2^{k-1}\) residue classes of \(x\pmod{2^k}\) admit a corresponding \(y\).

This immediately shows why simply increasing 2-adic precision does not isolate the factor coordinate: half of the \(x\)-classes survive for every \(k\).

\subsection{The involution}

For fixed \(y\), the two roots of \(P(x,y)\) satisfy \eqref{eq:vieta}. Thus the modular solution set has the exact involution
\begin{equation}
\boxed{(x,y)\longleftrightarrow(1-y-x,y).}
\label{eq:involution}
\end{equation}
Experiments verified zero involution failures for powers \(2^k\) over multiple precisions. The true pair \(x_p,x_q\) always forms one of these two-cycles.

\subsection{Fixed points and 2-adic square roots}

A fixed point of \eqref{eq:involution} satisfies
\[
x=1-y-x,
\]
so
\[
y=1-2x.
\]
Substitution into \(P\) gives the exact identity
\begin{equation}
\boxed{
P(x,1-2x)=N-(x-s-1)^2.
}
\label{eq:fixedsquare}
\end{equation}
Putting
\[
r=x-s-1,
\]
we obtain
\begin{equation}
\boxed{r^2\equiv N\pmod{2^k}.}
\label{eq:2roots}
\end{equation}
Thus the fixed points are precisely shifted square roots of \(N\) modulo \(2^k\).

For odd \(N\), these roots exist exactly when \(N\equiv1\pmod8\), in which case there are four roots modulo \(2^k\) for \(k\ge3\). Experiments confirmed the direct correspondence between the fixed points of \(P\) and the modular square roots of \(N\).

\section{The square-root branches and the factor gap}
\label{sec:sqrtbranches}

Let \(r\) be any square root of \(N\) modulo \(2^k\):
\begin{equation}
 r^2\equiv N=pq\pmod{2^k}.
\end{equation}
Then
\begin{align}
(r-p)(r+p)&=r^2-p^2\\
&\equiv pq-p^2\\
&=p(q-p)\pmod{2^k}.
\end{align}
Hence
\begin{equation}
\boxed{(r-p)(r+p)\equiv p(q-p)\pmod{2^k}.}
\label{eq:rbranchp}
\end{equation}
Similarly,
\begin{equation}
\boxed{(r-q)(r+q)\equiv q(p-q)\pmod{2^k}.}
\label{eq:rbranchq}
\end{equation}
Since \(p,q,r\) are odd, both factors on the left have controlled 2-adic valuations. Up to truncation at the modulus, one obtains the exact valuation conservation law
\begin{equation}
\boxed{
\vp(r-p)+\vp(r+p)=\vp(q-p),
}
\label{eq:v2gap}
\end{equation}
for the appropriate branch, with the symmetric relation obtained by interchanging \(p\) and \(q\).

Experiments found the branch valuations to be permutations of \((1,t-1)\), where \(t=\vp(q-p)\), until the modulus cap is reached. This is a genuine and exact 2-adic structure.

However, targeted collision experiments showed that \(\vp(q-p)\) is not determined by low-bit fingerprints such as
\[
(N\bmod 2^k,\ s\bmod2^k)
\]
even for moderately large \(k\). Explicit collisions with different values of \(\vp(q-p)\) survived targeted searches at \(k=13\) and \(k=14\) across several prime sizes. Therefore the valuation structure cannot be reduced to a simple low-bit fingerprint of \((N,s)\).

\section{Quadratic-discriminant information modulo \(2^k\)}
\label{sec:discmod}

The discriminant relation \eqref{eq:disc} provides a second modular viewpoint. For the true \(y\),
\[
\Delta(y)=(q-p)^2.
\]
Consequently the square residue \(\Delta(y)\bmod2^k\) determines the valuation of \(q-p\) exactly whenever the square is nonzero modulo the chosen precision. When the square vanishes modulo \(2^k\), only the lower bound \(2\vp(q-p)\ge k\) is visible.

This is standard 2-adic square information rather than a factorization shortcut. In finite experiments, increasing \(k\) quickly makes the valuation exact for most cases, but the same modular data does not determine the actual factor pair.

\section{Quadratic-residue sieving of the Fermat coordinate}
\label{sec:sieve}

Although the \(y\)-search is exactly Fermat, modular information can still accelerate the search. From \eqref{eq:disc}, any valid factor coordinate must satisfy
\begin{equation}
\Delta(y)=(2s+1+y)^2-4N
\end{equation}
is a quadratic residue modulo every modulus.

For a modulus \(m\), define
\begin{equation}
\mathcal S_m=\{u^2\bmod m:u\in\Z\}.
\end{equation}
A candidate \(y\) survives the corresponding sieve if
\begin{equation}
\Delta(y)\bmod m\in\mathcal S_m.
\end{equation}
The true \(y\) necessarily survives every such sieve.

\subsection{Local densities}

For odd prime moduli, approximately half of the residue classes are quadratic residues. Powers of two are more structured. In the experiments, modulo \(8\) approximately one half of the candidates survived, while modulo \(16\) approximately one quarter survived. Combining several small moduli multiplies these local filtering effects to first approximation.

The following table summarizes the observed aggregate reduction in Experiment 58 for 100 random 18-bit-prime semiprimes. The raw number of tested \(y\)-values was 397110 across all cases.

\begin{center}
\begin{tabular}{lrrr}
\toprule
Moduli & Survivors & Reduction & Approx. factor \\
\midrule
\([8]\) & 198555 & 0.50000000 & 2.0\\
\([16]\) & 99305 & 0.74993075 & 4.0\\
\([16,3]\) & 48752 & 0.87723301 & 8.1\\
\([16,3,5]\) & 25147 & 0.93667498 & 15.8\\
\([16,3,5,7]\) & 12439 & 0.96867619 & 32.0\\
\([16,3,5,7,11]\) & 6099 & 0.98464154 & 65.1\\
\([16,3,5,7,11,13]\) & 3067 & 0.99227670 & 129.5\\
\([16,3,5,7,11,13,17]\) & 1590 & 0.99599607 & 249.8\\
\([16,3,5,7,11,13,17,19]\) & 812 & 0.99795523 & 489.0\\
\([16,3,5,7,11,13,17,19,23]\) & 450 & 0.99886681 & 882.5\\
\bottomrule
\end{tabular}
\end{center}

The final wheel therefore removes approximately \(99.8867\%\) of the raw candidates, leaving only about one candidate in 882 on average.

\subsection{True-coordinate survival}

In the same experiment, the true \(y\) survived the full modulus set in all 100 cases. This is of course algebraically necessary, but it serves as an implementation check:
\begin{equation}
\boxed{100/100\text{ true-}y\text{ survival}.}
\end{equation}

The number of sieve survivors before the true \(y\) was small: the observed rank ranged from 1 to 22, with mean 4.50.

\subsection{Does the sieve know the true }\(y\)?

A very important distinction arises here. The final sieve density is approximately
\begin{equation}
\rho\approx\frac{450}{397110}\approx0.00113319.
\end{equation}
For a candidate interval containing \(y+1\) points, a uniform sieve of density \(\rho\) would predict an expected survivor rank of approximately
\begin{equation}
\rho(y+1).
\end{equation}
The observed mean rank of 4.50 is of exactly the same scale.

For example, a case with \(y=14073\) had rank 14, while
\[
14074\rho\approx15.95.
\]
Another case with \(y=16973\) had rank 16, compared with
\[
16974\rho\approx19.24.
\]
Conversely, very small true values such as \(y=147\) or \(y=29\) often produced rank 1 simply because too few candidates existed before the first survivor.

The present evidence therefore supports the conservative interpretation:
\begin{quote}
The quadratic-residue sieve is a strong candidate reduction mechanism, but the observed rank of the true \(y\) is quantitatively consistent with ordinary sieve density. No independent evidence has yet been established that the sieve preferentially selects the factor-producing \(y\).
\end{quote}

\section{Experiment log and exact results}
\label{sec:experiments}

This section records the principal computational experiments in a compact mathematical form.

\subsection*{Experiment 23: transformed quadratic \(H\)}
The arithmetic-sequence transform produced \(T(x)=3^{x-1}H(x)\) with \(H\) given by \eqref{eq:H}. The exact factor-coordinate evaluation \eqref{eq:Hxtrue} was verified. A single zero of \(H\) occurred in the tested data (\(N=6649=61\cdot109\), \(s=81\)), showing that algebraic degeneracies can occur but are not sufficient to isolate a factor.

\subsection*{Experiments 37--40: low-bit fingerprints and solvability}
Attempts to recover the factor coordinate from \((N\bmod2^k,s\bmod2^k)\) showed increasing apparent uniqueness at small sample sizes, but later exhaustive/targeted collision tests destroyed universal low-bit claims. Exact algebra established the true criterion \eqref{eq:parity}: solvability modulo \(2^k\) is equivalent to \(x\equiv s\pmod2\), leaving exactly half of all \(x\)-classes.

\subsection*{Experiments 41--43: 2-adic root sets}
Enumeration of modular solutions verified that the true coordinate survives at every tested precision. The solution set carries the exact involution \eqref{eq:involution}; its fixed points are exactly the shifted square roots \eqref{eq:2roots}.

\subsection*{Experiments 44--47: square-root valuations}
The relation between square-root branches and \(q-p\) was established exactly through \eqref{eq:rbranchp}--\eqref{eq:v2gap}. However, targeted collision tests showed that the valuation of the factor gap is not determined by low-bit fingerprints of \((N,s)\).

\subsection*{Experiment 48: discriminant identity}
The corrected discriminant formula \eqref{eq:disc} was verified for 200/200 cases. Recovering the exact factor gap from the discriminant square was likewise successful whenever 2-adic precision was sufficient to avoid the zero-square ambiguity.

\subsection*{Experiment 49: equivalence with Fermat}
The \(y\)-coordinate search matched Fermat factorization exactly. The mapping \(A=s+\ceil{y/2}\) introduces only a coordinate shift.

\subsection*{Experiment 50: geometry of \(H\)}
Across 5000 cases, all transformed quadratics had real roots, but the real roots, vertex, and local-minimum features did not locate \(x_p\). Simple predictors scored zero correct cases. This rules out the most obvious real-analytic interpretations of \(H\).

\subsection*{Experiment 51: exact diagonal parameterization}
The identities \eqref{eq:diaglinear}--\eqref{eq:slope} were tested over 200 random cases. The linearization, divisor equivalence, and diagonal slope identities had zero failures. The diagonal is therefore an exact divisor parameterization, not a heuristic approximation.

\subsection*{Experiments 52--53: transformed diagonal search}
The transformed diagonal sequence \(B_n(r)\) was checked against its closed form and against direct gcd detection. All identities held. The transform recovered known factors when \(r\) was already supplied, but searching \(r\) was equivalent to scanning divisors directly; no new selection principle emerged.

\subsection*{Experiments 55--56: \(a,b\) coordinates}
The equations \eqref{eq:yab} and \eqref{eq:abproduct} were verified on thousands of random semiprimes. Searching \(y\) and enforcing integral reconstruction of \(a,b\) remained equivalent to a Fermat square test.

\subsection*{Experiment 57: small-modulus discriminant sieve}
The sieve based on quadratic residues of \(\Delta(y)\) preserved the true \(y\) in every case. Combining \(16,3,5,7,11,13\) reduced the full candidate workload by about \(99.3\%\), from 334556 raw candidates to approximately 2323 survivors in the tested sample.

\subsection*{Experiment 58: extended wheel}
Using \([16,3,5,7,11,13,17,19,23]\), 100 random cases produced 397110 raw candidates and only 450 survivors, a reduction of \(99.8867\%\). The true \(y\) survived in all 100 cases, with sieve rank between 1 and 22 and mean 4.50. This is consistent with the global survivor density and therefore does not, by itself, imply additional factor localization.

\section{What is exact, what is classical, and what remains open}
\label{sec:status}

It is useful to separate the findings into three categories.

\subsection{Exact structural results}
The following identities are exact consequences of the coordinate construction:
\begin{enumerate}[label=\arabic*.]
\item The weighted binomial identities \eqref{eq:sum0}--\eqref{eq:sum2}.
\item The polynomial identity \eqref{eq:Pfactor}.
\item Diagonal linearization \eqref{eq:diaglinear}.
\item Divisor equivalence \eqref{eq:xfromr}.
\item Exact diagonal slope \eqref{eq:slope}.
\item The transformed quadratic identity \eqref{eq:T-H}.
\item The transformed diagonal difference identity \eqref{eq:Cdifference}.
\item The discriminant identity \eqref{eq:disc} and its factor-gap specialization \eqref{eq:gap}.
\item The \(a,b\) equations \eqref{eq:diffab} and \eqref{eq:abproduct}.
\item The exact 2-adic solvability criterion \eqref{eq:parity}.
\item The involution \eqref{eq:involution} and its fixed-point square-root representation \eqref{eq:2roots}.
\item The valuation relations \eqref{eq:rbranchp}--\eqref{eq:v2gap}.
\end{enumerate}

\subsection{Exact but classically equivalent factorization procedures}
The following do not presently constitute new factorization algorithms:
\begin{enumerate}[label=\arabic*.]
\item Searching \(y\) and testing whether \(\Delta(y)\) is a square: this is Fermat factorization.
\item Solving for \(a,b\) from \(y\): this reduces to the same quadratic discriminant.
\item Searching the diagonal parameter \(r\) and using the transformed difference: this reduces to divisor testing or \(\gcd(r,N)\).
\item Studying the ordinary roots or vertex of \(H\): the experiments do not indicate a useful factor-location rule.
\end{enumerate}

\subsection{Potentially useful computational improvement}
The quadratic-residue sieve is the clearest algorithmic improvement discovered in this line of work. It does not change the mathematics of Fermat factorization, but it can reduce the number of expensive exact-square tests by orders of magnitude. For the tested 18-bit-prime cases and the modulus set \([16,3,5,7,11,13,17,19,23]\), the reduction was about \(885\times\).

The open question is whether the sieve can be organized into a fast wheel whose construction cost is negligible compared with the resulting savings, and whether larger or specially chosen moduli produce a nontrivial improvement beyond the approximately multiplicative density model.

\section{Further directions suggested by the structure}
\label{sec:future}

The present work narrows the search space for a genuinely new idea. Several directions are especially natural.

\subsection{Normalized rank testing}
For a case with true coordinate \(y\), define the observed survivor rank \(R(y)\) and compare
\begin{equation}
Q=\frac{R(y)}{(y+1)\rho},
\end{equation}
where \(\rho\) is the empirically calibrated sieve density. If \(Q\) is stable around one, this strongly supports the conclusion that the sieve carries no extra directional information.

\subsection{Structured wheel generation}
The full least-common-modulus construction can become unwieldy. A practical implementation can instead construct a moderate base wheel, for example using
\[
16\cdot3\cdot5\cdot7\cdot11\cdot13=240240,
\]
then apply additional modular filters for \(17,19,23,\ldots\) to the base residue classes. This avoids enumerating all residues of a huge product modulus.

\subsection{Beyond quadratic residues}
The current sieve uses only the condition that \(\Delta(y)\) be a square modulo \(m\). One may ask whether stronger local conditions can be extracted from the full polynomial \(P\), its derivative, or simultaneous congruences in \((x,y)\). Any such condition must be benchmarked against its random-residue density; otherwise an apparent reduction may simply be a standard local sieve.

\subsection{Cross-modulus correlations}
Although individual quadratic-residue conditions are largely local, the same \(y\) is evaluated against all moduli simultaneously. Measuring correlations between residue classes could reveal whether certain moduli are unusually informative when combined. This should be tested statistically rather than inferred from a small sample.

\subsection{Factor-gap information}
The relation
\[
\Delta(y)=(q-p)^2
\]
means that every surviving candidate implicitly predicts a candidate gap \(g=\sqrt{\Delta(y)}\). A future route would need a constraint on \(g\) that is stronger than local quadratic residuosity but does not amount to simply finding the square. The current 2-adic experiments provide structural information about \(\vp(g)\), but not a sufficient selector.

\section{Conclusion}
\label{sec:conclusion}

The arithmetic-sequence construction developed here leads to a surprisingly rigid algebraic picture. The weighted binomial identity transforms the original quadratic polynomial into a simpler quadratic \(H(x)\); diagonalization removes the quadratic term entirely and exposes a parameter \(r=s+z\) whose divisibility by \(N\) is equivalent to an integer zero on that diagonal. The normalized transformed difference recovers this same parameter exactly.

At the same time, the \(y\)-coordinate reveals that the entire polynomial system is a shifted representation of Fermat factorization. The discriminant is exactly \((q-p)^2\), and the alternative \((a,b)\) variables reduce to the same square condition. This prevents overinterpretation of the algebra: several elegant identities describe the factorization problem without making it easier in a complexity-theoretic sense.

The modular analysis is richer. Modulo \(2^k\), solvability has an exact parity criterion, the solution set carries an involution, its fixed points are shifted square roots of \(N\), and those square-root branches encode the 2-adic valuation of the factor gap. Yet targeted collision experiments rule out the simplest hypothesis that this information is determined by low-bit fingerprints of \(N\) and \(s\).

The strongest computational result is the quadratic-residue sieve on the Fermat coordinate. A carefully chosen collection of small moduli can eliminate roughly three orders of magnitude of candidate values while preserving the true factor coordinate. The current evidence, however, indicates that the true coordinate is merely an ordinary survivor of this sieve, not a statistically privileged one.

Thus the research program ends at a useful boundary. The arithmetic-sequence identities, diagonal parameterization, Fermat equivalence, 2-adic involution, and quadratic-residue wheel form a coherent exact framework. What remains is to find a condition that is simultaneously stronger than ordinary local residuosity, cheaper than an exact square test, and genuinely correlated with the hidden factor location. That is the point at which the present algebraic reparameterizations would have to give way to a new factor-discovery mechanism.

\appendix
\section{Representative computational verification}
\label{app:verification}

The computational experiments used random distinct odd primes and verified identities over the resulting semiprimes. A representative pseudocode skeleton is:

\begin{lstlisting}
def generate_case(bits):
    p = random_odd_prime(bits)
    q = random_odd_prime(bits)
    while q == p:
        q = random_odd_prime(bits)
    if p > q:
        p, q = q, p
    N = p*q
    s = isqrt(N)
    D = N - s*s
    x = s + 1 - p
    y = p + q - 2*s - 1
    return p, q, N, s, D, x, y


def P(N, s, D, x, y):
    return x*x + (y-1)*x + (D-s) - (s+1)*y


def discriminant(N, s, y):
    return (2*s + 1 + y)**2 - 4*N


def survives_residue_sieve(N, s, y, square_tables):
    delta = discriminant(N, s, y)
    for modulus, table in square_tables.items():
        if table[delta % modulus] == 0:
            return False
    return True
\end{lstlisting}

For reproducibility, each experiment script printed an explicit beginning and ending marker of the form
\begin{lstlisting}
START EXPERIMENT N
...
FINISHED EXPERIMENT N
\end{lstlisting}
so that the output streams could be compared unambiguously.

\section{Notation summary}
\begin{longtable}{p{0.20\textwidth} p{0.65\textwidth}}
\toprule
Symbol & Meaning\\
\midrule
$N$ & $pq$, the odd semiprime\\
$s$ & $\floor{\sqrt N}$\\
$D$ & $N-s^2$\\
$x$ & $s+1-p$\\
$y$ & $p+q-2s-1$\\
$P(x,y)$ & $x^2+(y-1)x+(D-s)-(s+1)y$\\
$z$ & $x+y$\\
$r$ & $s+z$, the diagonal divisor parameter\\
$H(x)$ & $4x^2-(s+4)x+3D-3s$\\
$A$ & $s+\ceil{y/2}$, Fermat's square-root coordinate\\
$a$ & $s-p$\\
$b$ & $q-s$\\
$v_2$ & 2-adic valuation\\
\bottomrule
\end{longtable}

\end{document}
